%%%%%%%%%%%%%%%%%%%%%%%%%%%%%%%%%%%%%%%%%%%%%%%%%%%%%
%			HLAVIČKA								%
%%%%%%%%%%%%%%%%%%%%%%%%%%%%%%%%%%%%%%%%%%%%%%%%%%%%%
\documentclass[openany,12pt]{memoir}

\usepackage[utf8]{inputenc} 
\usepackage[czech]{babel}
\usepackage[T1]{fontenc}
\usepackage[top=1.5cm, bottom=2cm, left=2cm, right=2cm]{geometry}  % --> NASTAVENÍ OKRAJŮ
\usepackage{fancyhdr}
\usepackage{graphicx}
\usepackage{lmodern}
\usepackage{xwatermark}
\usepackage{xcolor}
\usepackage{changepage}
\usepackage{newunicodechar}



%%%%%% Package na zpěvník
\usepackage[full]{leadsheets}%http://mirrors.nic.cz/tex-archive/macros/latex/contrib/leadsheets/leadsheets_en.pdf   --> dokumentace	
\definesongtitletemplate{empty}{} 
\setchords{
format = \bfseries,   %tučné akordy
minor = {mi},% 
input-notation = {german},%
output-notation = {german}%
}
\definesongtitletemplate{empty}{} 

\newlength{\drop}
\newwatermark[allpages,color=red!50,angle=0,scale=2, xpos=0,ypos=0]{\includegraphics[width=5cm]{obr/pozadi.jpg}} %--> dvojka na pozadí


%%%% Vlastní příkazy
\newcounter{Slokočet}   %Automatické číslování slok
\newcommand{\mezera}{\vspace*{0.5cm}}   %Horizontální odsazení slok
\newcommand{\stred}{5.2cm}   %%% Na zarovnání slok doprostřed, pozn. automatičtější zarovnávání na střed nejde
\newcommand{\refren}{\mezera \noindent \textbf{R:} } %refrén
\newcommand{\sloka}{\mezera \noindent \addtocounter{Slokočet}{1} \arabic{Slokočet}. } 	%sloka, která se automaticky čísluje
\newcommand{\ssloka}{\mezera \noindent}  % vlastní číslo sloky
\newcommand{\carka}{,\:}
\newcommand{\m}[1]{\color{white}{#1}}  %Pro akordy

%%%%%%%%%%%%%%%%%%%%%%%%%%%%%%%%%%%%%%%%%%%%%%%%%%%%%%
%			NÁVOD									 %
%%%%%%%%%%%%%%%%%%%%%%%%%%%%%%%%%%%%%%%%%%%%%%%%%%%%%%
%1. Věci v hlavičce IGNOROVAT
%2. Píseň psát do prostoru mezi \begin{song} a \end{song}
%3. další řádek se značí dvěma odsazeníma (= dvakrát stisknout enter)
%4. \refren vždy na začátku refrenu a \sloka na začátek sloky (automaticky se čísluje)
%  \ = alt gr + q ; [/] = alt gr f/g ; {/} = alt gr + b/n; ^ = alt gr + 3 + mezera
%Cokoliv napíšete do ^{  } se bude brát jako akord
%když se toto bude dotýkat nějakého slova (nebude mezi tím a slovem mezera)
%tak se akordy zjeví nad slovem, ale pište to před slovo
%Když se to nedotýká slova, tak akord lítá ve vzduchu a vytiskne se větší mezera
%První možnost je asi preferovanější
%5. Akordy stačí psát jen do první sloky, když se nezmění -- kytaristi to zvládnou



%<++>
%\usepackage{subfiles}
%</++>

\begin{document}

\pagestyle{empty}
%\addtocontents{toc}{\protect\thispagestyle{empty}}
%\tableofcontents \thispagestyle{empty}\newpage
\newgeometry{top=1.5cm, bottom = 0cm, left = 2cm, right = 2cm}



\begin{song}{title=\centering Hero of War \\\normalsize Rise against  \vspace*{-0.3cm}}  %% sem se napíše jméno songu a autor
\moveright \stred \vbox{      %Varianta č. 1  ---> Jeden sloupec zarovnaný na střed	

\sloka 
	He said, \uv{ ^{Em}Son,

	Have you seen the ^{H}world?

	Well, what would you ^{A}say
                   
	If I said that you ^{Emi}could?
               
	Just carry this ^{A}gun 
                   
	and you’ll even get ^{Emi}paid.}
                            
	I said, \uv{That sounds pretty ^{H}good.}
              
	Black leather ^{Emi}boots
              
	Spit-shined so ^{H}bright
                
	They cut off my ^{A}hair 
             
	but it looked ^{Emi}alright
                 
	We marched and we ^{A}sang
              
	We all became ^{Emi}friends
                     
	As we learned how to ^{H}fight


\refren
	A hero of ^{F\# mi}war

	Yeah that’s what I’ll ^{A}be
               
	And when I come ^{Emi}home
                        
	They’ll be damn proud of ^{H}me
               
	I’ll carry this ^{F\# mi}flag
                 
	To the grave if I ^{A}must
                        
	Because it’s flag that I ^{Emi}love
                  
	And a flag that I ^{H}trust


\sloka
	I kicked in the ^{Emi}door
             
	I yelled my ^{H}commands
                   
	The children, they ^{A}cried
              
	But I got my ^{Emi}man
            
	We took him ^{A}away
                
	A bag over his ^{Emi}face
                       
	From his family and his ^{H}friends

	They took off his ^{Emi}clothes
                 
	They pissed in ^{H}his hands
                
	I told them to ^{A}stop
              
	But then I ^{Emi}joined in
                 
	We beat him with ^{A}guns
                  
	And batons not just ^{Emi}once
               
	But again and ^{H}again


\refren


\sloka
	She ^{Emi}walked 
                    
	through bullets and ^{H}haze
                
	I asked her to ^{A}stop
               
	I begged her to ^{Emi}stay
               
	But she pressed ^{A}on
                
	So I lifted my ^{Emi}gun
             
	And I fired ^{H}away
    
	The ^{Emi}shells 
                    
	jumped through the ^{H}smoke
             
	And into the ^{A}sand
                        
	That the blood now had ^{Emi}soaked
        
	She ^{A}collapsed 
                   
	with a flag in her ^{Emi}hand
                
	A flag white as ^{H}snow


\sloka
	A hero of ^{F\# mi}war
                 
	Is that what the ^{A}see
                
	Just medals and ^{E}scars
                
	So damn proud of ^{H}me
                        
	And I brought home that ^{F\# mi}flag
               
	Now it gathers ^{A}dust
                        
	But it’s a flag that I ^{E}love
                       
	It’s the only flag I ^{H}trust


\refren
He said, \uv{ ^{Emi}Son, 
                   
have you seen the ^{H}world? 
                    
Well what would you ^{A}say, 
                   
if I said that you ^{Emi}could?}







\setcounter{Slokočet}{0}
\end{song}


\newpage


\end{document}
